\subsection*{Objectives}
This project develops a software-centric control-supervision module for a lower-limb robotic exoskeleton intended to assist heavy load transportation tasks ($>20$~kg) in industrial, logistics, and construction settings. The focus of the submitted work is real-time locomotion-intent recognition from inertial measurement unit (IMU) data, stable assist gating, and reproducible MATLAB-based evaluation on public datasets.

\subsection*{Methods}
The validated MATLAB baseline uses overlapping IMU windows of 100 samples at 100~Hz with a 50-sample step. Kalman fusion (\texttt{imufilter}) is included for simulation-side kinematic visualization rather than as mandatory classifier preprocessing.

The primary deployment classifier is a binary RBF support vector machine trained on a 30-dimensional handcrafted feature vector: five statistics per IMU slot across six slots. The final pipeline is centered on HuGaDB and uses a protocol-aware streaming policy with quality gating and official gyro-corruption handling. USC-HAD is retained as a separate benchmark dataset under the same feature layout. Binary labels group activities into active movement versus inactive movement for assist gating.

For analysis beyond binary gating, the repository also includes dataset-native multiclass ECOC models and optional sequence-based LSTM baselines. Binary and multiclass LSTM models operate on feature-by-time IMU sequences and are evaluated separately from the SVM baselines. Assist decisions in replay are stabilized by a finite-state machine with asymmetric hysteresis: three consecutive active predictions to engage assistance and five consecutive inactive predictions to release it.

\subsection*{Results}
Binary SVM five-fold out-of-fold accuracy from the committed evaluation outputs is \textbf{\BinaryOofAccUsc\%} on USC-HAD and \textbf{\BinaryOofAccHuGa\%} on HuGaDB under the multi-activity streaming policy. Native multiclass SVM accuracies on stratified subsamples are \textbf{\MulticlassOofAccUsc\%} on USC-HAD and \textbf{\MulticlassOofAccHuGa\%} on HuGaDB.

\ifnum\HasBinaryLstmMetrics=1
Binary LSTM holdout accuracy is \textbf{\LstmBinaryHoldoutAccUsc\%} on USC-HAD and \textbf{\LstmBinaryHoldoutAccHuGa\%} on HuGaDB.
\else
Binary LSTM metrics are optional and are exported after running the dataset-specific LSTM evaluation scripts.
\fi
\ifnum\HasMulticlassLstmMetrics=1
Native multiclass LSTM holdout accuracy is \textbf{\LstmMulticlassHoldoutAccUsc\%} on USC-HAD and \textbf{\LstmMulticlassHoldoutAccHuGa\%} on HuGaDB.
\else
Multiclass LSTM metrics are optional and are exported after running the dataset-specific sequence-model evaluation scripts.
\fi

Pseudo-real-time replay on held-out HuGaDB sessions shows that the combined classifier and FSM pipeline produces a stable binary assist command alongside IMU magnitude for visualization context. This report \textbf{emphasizes} the \textbf{multiclass LSTM} pseudo-real-time replay on HuGaDB subject~\ReplayHuGaDBSubject{}, session~\ReplayHuGaDBSession{} as the principal sequence-model demonstration: native twelve-class predictions over time on the same held-out session used for the binary replay figures.

\subsection*{Conclusions}
The project delivers a reproducible HuGaDB-centered MATLAB baseline for exoskeleton assist gating, with traceable scripts, saved metrics, evaluation figures, and replay outputs. The strongest validated \emph{classification} headline remains the binary SVM under the HuGaDB streaming policy, while the multiclass LSTM replay (subject~\ReplayHuGaDBSubject{}, session~\ReplayHuGaDBSession{}) is highlighted as the end-to-end sequence-model visualization for fine-grained activity behavior. Multiclass SVM and additional LSTM holdout metrics remain supporting analyses alongside that emphasis. The resulting software pipeline provides a solid foundation for future hardware integration, embedded deployment, and controlled experimental studies.
